% ===========================================================
%  SNA 期末報告 Proposal (修訂版 v2)
%  資料：TIPD/TIMD 2013-2022 多年份多維表 (半年期 + 年期)
%  Compile with XeLaTeX
% ===========================================================
\documentclass[11pt,a4paper]{article}

\usepackage[margin=2cm,top=1.8cm,bottom=1.8cm]{geometry}
\usepackage{xeCJK}
\setCJKmainfont{Noto Serif CJK TC}[BoldFont=Noto Serif CJK TC,ItalicFont=Noto Serif CJK TC]
\setCJKsansfont{Noto Sans CJK TC}
\setmainfont{TeX Gyre Termes}

\usepackage{titlesec}
\usepackage{enumitem}
\usepackage{xcolor}
\usepackage{parskip}
\usepackage{booktabs}
\usepackage{tabularx}
\usepackage{array}

\definecolor{secblue}{HTML}{1F4E79}
\definecolor{accentgreen}{HTML}{2E7D32}

\titleformat{\section}{\bfseries\large\color{secblue}}{\thesection.}{0.4em}{}
\titleformat{\subsection}{\bfseries\normalsize\color{accentgreen}}{\thesubsection}{0.4em}{}
\titlespacing*{\section}{0pt}{10pt}{4pt}
\titlespacing*{\subsection}{0pt}{6pt}{2pt}

\setlist[itemize]{leftmargin=1.4em, itemsep=2pt, topsep=2pt, parsep=0pt}
\setlength{\parskip}{3pt}
\linespread{1.10}

\newcolumntype{Y}{>{\raggedright\arraybackslash}X}

\title{}
\author{}
\date{}

\begin{document}

% ===== Title block =====
\begin{center}
  {\LARGE\bfseries 裂解與重組：以社會網絡分析探索泰雅族遷徙網絡的多尺度結構變遷}\\[4pt]
  {\small\itshape Fragmentation and Reconfiguration: Multi-scale Structural Transformation of\\Atayal Migration Networks through Social Network Analysis}\\[6pt]
  {\small 賴亭穎 \textbar{} 黃凡嘉 \quad\textbar\quad 國立臺灣師範大學 社會網絡分析 \quad\textbar\quad 2026 春季}
\end{center}
\vspace{-4pt}\hrule\vspace{4pt}

% =============================================================
\section{研究背景與問題}
泰雅族橫跨新北、桃園、新竹、苗栗、臺中、南投、宜蘭山區，是臺灣原住民族中分布領域最廣者。族群內分賽考列克與澤敖列兩大亞群，從少數祖源節點向外輻射，形成「共同祖源、多路線擴散」的網絡拓撲。其遷徙史可分為三個結構性斷裂期：(1)~\textbf{自發擴散期}（17 世紀末--1895）以血緣親族為單位、沿溪流建立散村，部落間透過 \textit{mulaxen galang} 聯盟結盟；(2)~\textbf{殖民強制移住期}（1895--1945）日治理蕃政策由上而下重組空間，戰後延續至 1971 年共遷移 46 村約 15{,}978 人；(3)~\textbf{戰後都市化期}（1945--）部落連結由地理鄰近轉向跨區域親族與同鄉網絡。

岡松參太郎指出泰雅族頭目實力弱、部落間爭執常以出草解決；當代田野中亦有族人表示「泰雅族是各原住民族中部落間互動最少的族群」。本研究將此田野觀察轉化為三個可檢驗的問題：

\vspace{-2pt}
\begin{itemize}
\item \textbf{Q1.} 殖民與都市化是「裂解」傳統網絡，抑或「重組」為新形式？
\item \textbf{Q2.} 「泰雅族部落間互動最少」之觀察，在哪個空間尺度與時間切片上成立？
\item \textbf{Q3.} 既有歷史敘事（廖守臣 1998、藤井志津枝 1997）所主張的網絡轉變，能否在量化網絡指標上被觀察到？
\end{itemize}

\subsection*{方法論貢獻}
Mills 等 (2013 \emph{PNAS}; 2018 \emph{Antiquity}) 以動態 SNA 分析北美原住民遷徙，但臺灣原住民遷徙研究至今未有系統性 SNA 分析。本研究結合 \textbf{歷史民族誌資料}（移川子之藏等 1935、廖守臣 1998）與\textbf{近十年 TIPD 多年份開放資料}，於三個空間尺度上建構互補的網絡，補上此一方法論空白。

% =============================================================
\section{資料來源}
\textbf{主要資料：TIPD/TIMD 多維表（中央研究院人文社會科學研究中心 林季平博士主持）}
\vspace{-4pt}
\begin{itemize}
\item 時間範圍：\textbf{2013--2022 共 10 年}，分為半年期（18 時段）與年期（9 時段）兩種粒度。
\item 每時段含三類子資料：\texttt{\_Intact}（存活人口）、\texttt{\_Decrease}（死亡／跨國移出）、\texttt{\_Increase}（出生／跨國移入）。
\item 篩 \texttt{EthnicityCate=2} 取得泰雅族子集；2022 年泰雅族總樣本約 33{,}762 人。
\item 變數涵蓋出生地（縣市）、現居地（鄉鎮市區）、五歲年齡組、性別、教育程度、婚姻狀態。
\end{itemize}

\textbf{補充資料：}
\vspace{-4pt}
\begin{itemize}
\item \textbf{歷史民族誌}：移川子之藏等《臺灣高砂族所屬系統之研究》(1935) 第二卷系譜、廖守臣《泰雅族的文化：部落遷徙與拓展》(1998)、藤井志津枝《理蕃》(1997)、臺灣總督府《理蕃誌稿》。
\item \textbf{TICD}（臺灣原住民部落基礎資料）：部落屬性、座標、所屬亞群。
\item \textbf{GIS}：內政部國土測繪中心鄉鎮市區行政界 shapefile。
\end{itemize}

% =============================================================
\section{研究方法：三層網絡設計}

本研究於三個空間尺度建構互補網絡：

\begin{center}\small
\renewcommand{\arraystretch}{1.3}
\begin{tabularx}{\textwidth}{|l|Y|Y|Y|}
\hline
\textbf{層} & \textbf{Layer C：歷史部落} & \textbf{Layer A：縣市時序} & \textbf{Layer B：鄉鎮共源} \\
\hline
\textbf{時間} & 17 世紀末--1945 & 2013--2022（18/9 時段）& 2022 快照 \\
\textbf{節點} & 歷史蕃社（300+）& 22 縣市 & 170 鄉鎮（人口 $\geq$ 30）\\
\textbf{邊} & 有向：系譜記載之移住事件 & 有向加權：跨縣市終身遷徙流 & 無向加權：居民出生地分布之相似度 \\
\textbf{邊權重} & 遷徙戶／人數 & 該時段人口數 & Brainerd--Robinson 相似度 \\
\textbf{資料源} & 移川 1935、廖守臣 1998、《理蕃誌稿》 & TIMD 2013--2022 & TIMD 2022 \\
\hline
\end{tabularx}
\end{center}

\subsection{Layer A 時序網絡建構（重點創新）}
利用 TIMD 多年份特性，建構真正的\textbf{時序縣市遷徙網絡}。對每一時段 $t$（2013H1, 2013H2, \ldots, 2022H2）：
\begin{itemize}
\item 節點：22 個縣市（升格代碼正規化後）；
\item 邊權重 $w_{ij}^{(t)}$：時段 $t$ 中「出生於縣市 $i$ 且現居縣市 $j$」之人數；
\item 同時依\textbf{五歲年齡組}建構 cohort 子網絡，追蹤同一出生世代隨時間的網絡變化。
\end{itemize}

此設計使本研究得以\textbf{分離 cohort effect 與 age effect}——觀察性人口學研究最棘手的混淆問題：
\begin{itemize}
\item \textbf{Age effect}：同一個 1957 出生世代在 2013（56 歲）、2018（61 歲）、2022（65 歲）之網絡演化；
\item \textbf{Cohort effect}：1988 世代在 2013（25 歲）與 1997 世代在 2022（25 歲）之網絡比較。
\end{itemize}

\subsection{Layer B 鄉鎮共源相似度網絡}
\textbf{譯自 Mills 等 (2013) PNAS 之方法}：以鄉鎮居民「出生縣市分布」為節點屬性向量，邊權重為 Brainerd--Robinson 相似度（介於 0--1），取上 10\% 強邊建網絡。此設計揭示「具有相同人口磁場」之鄉鎮叢——預期山地原鄉與其下游都會延伸應屬同社群。

\subsection{Layer C 歷史部落網絡}
以移川 1935 第二卷之系譜記錄與廖守臣 1998 之遷徙考證重建 pre-1945 部落級網絡。每筆「某氏族從 A 社遷至 B 社」轉化為一條有向邊。藉宇野圓空 (1935) 書評之批判性提醒：發祥傳說需區分歷史層與神話層。

\subsection{分析指標與穩健性檢驗}
\begin{itemize}
\item \textbf{中心性}：in/out-degree、betweenness（取 $\text{distance}=1/w$）、eigenvector、PageRank；
\item \textbf{社群偵測}：Louvain (Blondel et al. 2008) 為主，Leiden (Traag et al. 2019) 作穩健性對照（報告 ARI 與 NMI）；
\item \textbf{結構洞}：Burt (1992) effective size 與 constraint；
\item \textbf{弱連結／橋接}：Granovetter (1973) 之 local bridge 偵測（移邊後最短路徑 $\geq 2$）；
\item \textbf{Null model}：configuration model 隨機化 200 次 $\rightarrow$ modularity z-score；spatial null 檢驗空間--社會獨立性。
\end{itemize}

% =============================================================
\section{初步成果（pilot analysis on 2022 \_Intact）}

\subsection*{Layer A 縣市世代比較}
以 2022 單時點 \_Intact 進行 cohort 準時序分析（待擴展為真實時序）：
\begin{itemize}
\item 網絡密度由 C4（65+ 歲）之 0.36 升至 C2（25--44 歲）之 0.70；
\item Modularity Q 對 configuration null 之 z-score：C4 = +3.1（顯著有社群）$\to$ C3 = $-$2.5（顯著無社群）—— 揭示\textbf{社群結構崩解之關鍵世代為 C3（1957--1977 出生）}；
\item 桃園、新北在 C1 世代（$<$25 歲）出現淨流出（$-754$、$-576$），呈現「\textbf{二代離散}」——當代都會原鄉的下一輪擴散。
\end{itemize}

\subsection*{Layer B 鄉鎮共源網絡}
170 個鄉鎮、Brainerd--Robinson 上 10\% 強邊：
\begin{itemize}
\item \textbf{Modularity Q = 0.654}（vs 隨機網絡約 0.01），偵測出 10 個地理--歷史社群；
\item 山地原鄉（尖石、五峰、大同、南澳、仁愛、和平等）與其下游都會（竹東、宜蘭市、埔里、北屯等）多歸於同一社群—\textbf{支持「流域廊道延續」假說}；
\item 橋接型鄉鎮：竹南、寶山、楊梅、三義、梧棲——均為縣界都會邊緣；
\item Leiden 演算法穩健性檢驗：所有網絡 ARI $\geq$ 0.84，分區結果穩健。
\end{itemize}

% =============================================================
\section{預期成果}
本研究預期回答提案開頭三個核心問題：
\begin{itemize}
\item \textbf{針對 Q1}：證明「裂解／重組」需於不同尺度分別回答——縣市層級顯示傳統聯盟邊界已被打散（modularity 降至 null 之下），但鄉鎮層級仍保留強流域廊道結構（Q = 0.654）。
\item \textbf{針對 Q2}：「部落互動最少」之觀察主要成立於縣市層級之 C3 世代之後；鄉鎮層級之社群結構仍清晰可辨。
\item \textbf{針對 Q3}：藉時序網絡資料分離 cohort effect 與 age effect，量化評估廖守臣與藤井志津枝之歷史敘事。
\end{itemize}

\subsection*{方法論貢獻}
\vspace{-2pt}
\begin{itemize}
\item 首次系統性將 SNA 應用於臺灣原住民族遷徙研究；
\item 整合 1930 年代田野系譜與 2020 年代開放資料的多時間尺度設計；
\item 三層網絡之多尺度設計可推廣至其他原住民族與少數族群之遷徙分析。
\end{itemize}

% =============================================================
\section{研究限制與倫理}
\vspace{-2pt}
\begin{itemize}
\item TIMD 鄉鎮層級之有向 OD 流尚未公開，目前僅以「出生地 $\to$ 現居地」推導終身遷徙；
\item 歷史層資料品質不均，殖民期文書需批判性使用；
\item 倫理：遵守 CARE 原則與 TIPD 使用規範；具體部落資訊以匿名或區域層級呈現；研究結果擬回饋至泰雅族社群諮詢意見。
\end{itemize}

% =============================================================
\vspace{4pt}\hrule\vspace{2pt}
{\footnotesize\textbf{核心參考文獻：}%
Granovetter (1973); Burt (1992); Freeman (1978); Watts \& Strogatz (1998); Barab\'asi \& Albert (1999); Blondel et al.\ (2008); Traag, Waltman \& van Eck (2019); Mills et al.\ (2013, 2018); Aref \& Wilson (2018); 移川子之藏等 (1935); 衛惠林 et al.\ (1984); 廖守臣 (1998); 藤井志津枝 (1997); 林季平 TIPD/TIMD (OSF: 6rpz9).}

\end{document}
